\chapter{Experimental Design}
\label{ch:experimental}

\section{Introduction}
\label{sec:exp-intro}

This chapter describes the experimental design used to evaluate the NirmalNode prototype introduced in Chapter~3: the physical setup and prototype build, the sensor calibration procedure, the data collection protocol, the evaluation metrics applied to the Adaptive AI module and to the purification unit, the estimated energy/cost profile of the system, and a discussion of the observed results and their limitations. Because the specific numerical results (accuracy figures, efficiency percentages, measured power draw) depend on hardware that only the author can run and measure, every quantitative result in this chapter is marked with a \TODO{} placeholder indicating exactly which measured value must be inserted once the corresponding experiment has been executed. The experimental \emph{procedure} described around each placeholder is, however, complete and ready to follow.

\section{Experimental Setup}
\label{sec:exp-setup}

\subsection{Prototype Description}
\label{sec:exp-prototype}

The bench-scale prototype consists of a single NirmalNode unit built around an ESP32 Dev Module, wired to a PMS5003 particulate sensor (UART), an MiCS-4514 dual-element gas sensor (analog, via the ESP32's onboard ADC), an MQ135 and MQ136 gas sensor pair (analog), an MP135 supplementary sensor (analog), and an optional BME280 environmental sensor (I$^2$C), as specified in Section~\ref{sec:meth-hardware}. The purification train -- cyclone separator, HEPA filter stage, and activated-carbon stage, drawn through by a 12~V DC fan -- is built as a compact, worker-facing enclosure sized to direct purified air into a single worker's immediate breathing zone rather than a whole room. Power for the prototype is supplied by a bench DC power supply (regulated 5~V for the ESP32/sensor stage and 12~V for the fan), consistent with the "prototype" power configuration specified in the project brief; a Li-Po battery pack is identified as the target power source for a field-deployable version (Section~\ref{sec:future-solar}).

\FIGPLACE{fig:prototype-photo}{Assembled bench-scale NirmalNode prototype (sensor stage and purification stage)}{0.8}
\FIGPLACE{fig:prototype-wiring}{Wiring/breadboard diagram of the ESP32 and connected sensors}{0.8}

\subsection{Test Environment}
\label{sec:exp-environment}

Two test environments are used, consistent with the escalation strategy recommended for prototype-stage embedded systems work:

\begin{enumerate}
    \item \textbf{Controlled laboratory environment.} The prototype is operated in a laboratory setting where a controlled pollutant source (e.g., incense smoke, a small combustion source, or a nebulized particulate source, subject to safety approval) can be introduced at a known distance from the sensor inlet, allowing the sensor readings and purifier response to be compared against a repeatable, human-controlled ground truth.
    \item \textbf{Representative industrial hotspot.} Where site access permits, the prototype is additionally deployed for a limited monitoring period at a real industrial hotspot (e.g., a welding bay or generator room) of the type described in Section~\ref{sec:industrial-pollution-bd}, to validate that the sensing and prediction pipeline behaves sensibly on real, non-laboratory data. \TODO{Name and describe the actual site(s) used, with any necessary permission/anonymization notes.}
\end{enumerate}

\section{Sensor Calibration}
\label{sec:exp-calibration}

Low-cost MOS gas sensors (MQ135, MQ136, MiCS-4514, MP135) exhibit sensor-to-sensor variation and a non-linear response to gas concentration, and therefore require calibration before their raw analog-to-digital converter (ADC) output can be treated as a meaningful concentration estimate. The calibration procedure followed is:

\begin{enumerate}
    \item \textbf{Burn-in.} Each MOS sensor is powered continuously for a burn-in period (typically 24--48 hours, per manufacturer datasheet guidance) before any calibration reading is taken, to allow the sensing element's resistance to stabilize.
    \item \textbf{Baseline (clean-air) calibration.} With the sensor exposed to outdoor or otherwise verified clean air, the sensor's baseline resistance ratio $R_0$ is recorded and used as the reference point for the manufacturer-supplied sensitivity curve.
    \item \textbf{Reference-instrument comparison.} Where available, sensor output is compared against a reference-grade or commercially calibrated air quality instrument over a range of concentrations, and a correction curve (e.g., linear or power-law fit between raw sensor output and reference concentration) is derived for each channel.
    \item \textbf{Temperature/humidity compensation.} Because MOS sensor resistance is sensitive to ambient temperature and humidity, BME280 readings are used to apply a compensation correction to the raw gas-sensor output, following the manufacturer-recommended compensation curves where available.
\end{enumerate}

\noindent \TODO{Insert the actual calibration coefficients / correction curves obtained for each sensor channel, and the reference instrument used, once calibration has been performed.}

\section{Data Collection}
\label{sec:exp-data-collection}

Data is logged at the sampling interval $\Delta t$ defined in Section~\ref{sec:meth-system-flow}, with each record consisting of the sensor readings listed in Section~\ref{sec:meth-hardware} together with timestamp, day-of-week, and location identifier fields, and stored in the cloud database described in Section~\ref{sec:meth-cloud}. Table~\ref{tab:exp-dataset} summarizes the composition of the dataset used to train and evaluate the Adaptive AI model.

\begin{table}[htbp]
\centering
\caption{Composition of the dataset used for model training and evaluation}
\label{tab:exp-dataset}
\begin{tabular}{@{}p{5cm}p{7cm}@{}}
\toprule
\textbf{Component} & \textbf{Description} \\
\midrule
Historical (pre-training) data & \TODO{Source and size of the historical Bangladesh air-quality dataset used for initial model training} \\
Laboratory-collected data & \TODO{Duration and number of records collected under the controlled test in Section~\ref{sec:exp-environment}} \\
Field (hotspot) data & \TODO{Duration and number of records collected at the real industrial hotspot, if applicable} \\
Sampling interval ($\Delta t$) & 60 seconds (prototype configuration) \\
Fields per record & PM1.0, PM2.5, PM10, CO, NO$_2$, VOC, NH$_3$, H$_2$S, smoke index, temperature, humidity, pressure, timestamp, day, location ID \\
Train / validation / test split & \TODO{Specify split strategy -- note that for an \emph{incremental} learner, a prequential (test-then-train) evaluation protocol is generally more appropriate than a fixed train/test split; see Section~\ref{sec:exp-metrics-procedure}} \\
\bottomrule
\end{tabular}
\end{table}

\section{Evaluation Metrics and Procedure}
\label{sec:exp-metrics-procedure}

\subsection{Prediction Accuracy}
The Adaptive AI model's forecasting accuracy is evaluated using the MAE, RMSE, and R$^2$ metrics defined in Equations~\ref{eq:mae}--\ref{eq:r2}, and, for the derived binary hazard-classification task (i.e., whether the forecast correctly identifies an upcoming hazardous period), using Precision, Recall, and F1-score as defined in Equation~\ref{eq:prf1}. Because the model is an \emph{incremental} learner rather than a model trained once on a fixed dataset, evaluation follows a \textbf{prequential (``test-then-train'') protocol}: for each new incoming record, the current model first produces its forecast $\hat{y}_{t+h}$ (which is scored against the eventual true value once time $t+h$ has elapsed), and only \emph{after} this evaluation step is the model updated with the new observation via Equation~\ref{eq:sgd} or~\ref{eq:pa} (or the corresponding tree/ensemble update). This protocol directly measures the model's real-world forecasting accuracy at each point in its deployment lifetime, and, critically, allows the accuracy trend over time to be plotted to demonstrate (or refute) the central Adaptive AI claim that accuracy improves as more site-specific data accumulates (Section~\ref{sec:exp-adaptivity-result}).

\subsection{Purification Efficiency}
Purification efficiency is measured as the percentage reduction in pollutant concentration between the unfiltered inlet air and the filtered outlet air of the purification unit, for each pollutant channel $p$:
\begin{equation}
\eta_p = \left(1 - \frac{C_{p,\text{out}}}{C_{p,\text{in}}}\right) \times 100\%
\label{eq:purification-efficiency}
\end{equation}
measured using two PMS5003/gas-sensor probes positioned at the inlet and outlet of the cyclone--HEPA--activated-carbon train respectively, under a controlled pollutant source (Section~\ref{sec:exp-environment}).

\subsection{Energy and Power Consumption}
System power draw is measured separately for (i) the sensing/compute stage (ESP32 + sensors, idle and active/transmitting) and (ii) the purification stage (fan, ON state), using an inline USB/DC power meter, to establish the total energy budget per duty cycle and to project battery life for the future Li-Po-powered field version (Section~\ref{sec:future-solar}).

\section{Results}
\label{sec:exp-results}

\subsection{Predictive Accuracy Results}
\label{sec:exp-accuracy-result}
\TODO{Insert MAE, RMSE, R\textsuperscript{2}, Precision, Recall, and F1-score results, ideally reported per pollutant channel (PM2.5, CO, etc.) and per candidate incremental model (SGD Regressor, Passive-Aggressive Regressor, Hoeffding Tree, Adaptive Random Forest), in a comparison table, together with the confusion matrix for the hazard-classification task.}

\begin{table}[htbp]
\centering
\caption{Predictive accuracy of candidate incremental-learning models (template -- to be completed with measured results)}
\label{tab:exp-model-comparison}
\begin{tabular}{@{}lccccc@{}}
\toprule
\textbf{Model} & \textbf{MAE} & \textbf{RMSE} & \textbf{R\textsuperscript{2}} & \textbf{F1-score} & \textbf{Update time/sample} \\
\midrule
SGD Regressor & \TODO{} & \TODO{} & \TODO{} & \TODO{} & \TODO{} \\
Passive-Aggressive Regressor & \TODO{} & \TODO{} & \TODO{} & \TODO{} & \TODO{} \\
Hoeffding Tree & \TODO{} & \TODO{} & \TODO{} & \TODO{} & \TODO{} \\
Adaptive Random Forest & \TODO{} & \TODO{} & \TODO{} & \TODO{} & \TODO{} \\
\bottomrule
\end{tabular}
\end{table}

\subsection{Adaptivity / Incremental-Learning Result}
\label{sec:exp-adaptivity-result}
To directly test the central Adaptive AI claim, prequential error (Section~\ref{sec:exp-metrics-procedure}) is plotted against time/sample index over the course of the deployment period, both (a) at a single fixed hotspot, to show whether accuracy improves as the model accumulates site-specific data, and (b), if a relocation test is performed, immediately before and after the unit is moved to a new location, to show whether and how quickly the model re-adapts. \TODO{Insert the resulting prequential-error-over-time plot(s) and a brief quantitative summary (e.g., \% reduction in rolling MAE from week 1 to week $n$).}

\FIGPLACE{fig:prequential-error-plot}{Prequential (test-then-train) error of the Adaptive AI model over the deployment period}{0.85}

\subsection{Purification Efficiency Results}
\label{sec:exp-purification-result}
\TODO{Insert measured $\eta_p$ (Equation~\ref{eq:purification-efficiency}) for each pollutant channel, along with the pollutant concentration time series at the purifier inlet and outlet during a representative test run.}

\begin{table}[htbp]
\centering
\caption{Measured purification efficiency by pollutant channel (template -- to be completed with measured results)}
\label{tab:exp-purification-efficiency}
\begin{tabular}{@{}lccc@{}}
\toprule
\textbf{Pollutant} & \textbf{Inlet conc.\ ($C_{p,\text{in}}$)} & \textbf{Outlet conc.\ ($C_{p,\text{out}}$)} & \textbf{Efficiency $\eta_p$ (\%)} \\
\midrule
PM2.5 & \TODO{} & \TODO{} & \TODO{} \\
PM10 & \TODO{} & \TODO{} & \TODO{} \\
CO & \TODO{} & \TODO{} & \TODO{} \\
NO$_2$ & \TODO{} & \TODO{} & \TODO{} \\
VOC & \TODO{} & \TODO{} & \TODO{} \\
NH$_3$ / H$_2$S & \TODO{} & \TODO{} & \TODO{} \\
\bottomrule
\end{tabular}
\end{table}

\subsection{Power Consumption Results}
\label{sec:exp-power-result}
\begin{table}[htbp]
\centering
\caption{Measured power consumption by subsystem (template -- to be completed with measured results)}
\label{tab:exp-power}
\begin{tabular}{@{}lcc@{}}
\toprule
\textbf{Subsystem} & \textbf{State} & \textbf{Power draw (W)} \\
\midrule
ESP32 + sensors & Idle / sensing & \TODO{} \\
ESP32 + sensors & Wi-Fi transmitting & \TODO{} \\
Purifier fan & OFF & 0 \\
Purifier fan & ON & \TODO{} \\
\textbf{Total (worst case, fan ON + transmitting)} & & \TODO{} \\
\bottomrule
\end{tabular}
\end{table}

\subsection{Cost Analysis}
\label{sec:exp-cost}
Table~\ref{tab:exp-cost} lists the bill of materials for the bench-scale prototype. Costs are indicative, prototype-quantity, retail-level estimates and should be re-verified against current local supplier pricing at the time of any deployment costing exercise, since component prices for imported electronics in Bangladesh fluctuate with exchange rates and import duty.

\begin{table}[htbp]
\centering
\caption{Indicative bill of materials for one NirmalNode prototype unit}
\label{tab:exp-cost}
\begin{tabular}{@{}p{5.5cm}p{3.5cm}p{4cm}@{}}
\toprule
\textbf{Component} & \textbf{Qty} & \textbf{Estimated unit cost} \\
\midrule
ESP32 Dev Module & 1 & \TODO{} \\
PMS5003 particulate sensor & 1 & \TODO{} \\
MiCS-4514 gas sensor & 1 & \TODO{} \\
MQ135 gas sensor & 1 & \TODO{} \\
MQ136 gas sensor & 1 & \TODO{} \\
MP135 sensor & 1 & \TODO{} \\
BME280 (optional) & 1 & \TODO{} \\
DC fan (12~V) & 1 & \TODO{} \\
Cyclone separator (fabricated/purchased) & 1 & \TODO{} \\
HEPA filter element & 1 & \TODO{} \\
Activated carbon filter element & 1 & \TODO{} \\
Enclosure, wiring, connectors, PCB/breadboard & 1 set & \TODO{} \\
\midrule
\textbf{Total (prototype, per unit)} & & \TODO{} \\
\bottomrule
\end{tabular}
\end{table}

\section{Discussion}
\label{sec:exp-discussion}

\TODO{Once the above results are available, discuss: (i) which incremental model(s) offered the best accuracy/update-time trade-off and why this is consistent (or not) with the theoretical properties described in Section~\ref{sec:meth-ai-pipeline}; (ii) whether the purification efficiency results are consistent with the HEPA-effectiveness literature reviewed in Section~\ref{sec:lr-purification}; (iii) whether the measured power/cost profile supports the Green IoT / low-cost-deployability claims of Section~\ref{sec:novelty}; and (iv) any unexpected sensor behaviour (cross-sensitivity, drift, calibration issues) observed during testing.}

\section{Limitations}
\label{sec:exp-limitations}

The evaluation presented in this chapter is subject to the following limitations, which should be read alongside the results above:

\begin{enumerate}
    \item \textbf{Single-unit, bench-scale prototype.} Results are obtained from a single physical unit rather than a multi-node deployment, so conclusions about network-level behaviour (e.g., across many simultaneous hotspots) are not directly supported by this evaluation.
    \item \textbf{Limited field-testing duration.} \TODO{State the actual duration of any real industrial-site testing performed, and note that the incremental model's long-run adaptivity claim is best supported by longer deployment periods than may have been feasible within the scope of this thesis.}
    \item \textbf{MOS gas sensor precision.} The MQ-series and MiCS-4514 sensors used are low-cost, broad-spectrum sensors with known cross-sensitivity and drift characteristics; their absolute concentration estimates are therefore less precise than a reference-grade electrochemical or optical gas analyzer, even after the calibration procedure of Section~\ref{sec:exp-calibration}.
    \item \textbf{Simulated vs.\ real hazard conditions.} Where a controlled laboratory pollutant source was used in place of an authentic uncontrolled industrial hazard, the resulting sensor and purifier response may not fully capture the variability of real industrial conditions.
    \item \textbf{Cost estimates are indicative.} The bill of materials in Table~\ref{tab:exp-cost} reflects prototype-quantity component pricing and does not account for potential economies of scale, import duty variation, or the cost of a production-grade (as opposed to bench-scale) enclosure and purification train.
\end{enumerate}

\section{Chapter Summary}
\label{sec:exp-summary}

This chapter has described the complete experimental design for evaluating NirmalNode: the prototype build and test environments, the sensor calibration procedure, the data collection protocol, the prequential evaluation procedure appropriate to an incremental learner, and the specific metrics (MAE, RMSE, R$^2$, Precision/Recall/F1, purification efficiency, power consumption, and cost) used to assess the system. All templates and procedures are in place; the corresponding result tables and figures are marked for completion once the physical experiments described here have been carried out. Chapter~5 concludes the thesis and outlines directions for future work.
